\subsection{可分多项式}

命题3. --- 设 \(f\) 是 \(\mathbf{K}[\mathbf{X}]\) 中的一个非零多项式，并令 \(\Omega\) 是 \(\mathbf{K}\) 的一个代数闭扩张。以下条件等价：

\begin{enumerate}
\def\labelenumi{\alph{enumi})}
\item
  多项式 \(f\) 与其导数 \(f'\) （在 \(\mathbf{K}[\mathbf{X}]\) 中）互素。
\item
  要么 \(\deg(f) = 0\) ，要么 \(\deg(f) > 0\) 且 \(\operatorname{dis}(f) \neq 0\) （IV，第～83页）。
\item
  存在一个扩张 \(\mathbf{L}\) （\(\mathbf{K}\) 的），使得 \(f\) 在 \(\mathbf{L}[\mathbf{X}]\) 中分解为若干个次数 \(\leqslant 1\) 的不同多项式的乘积。
\item
  \(f\) 在 \(\Omega\) 中的根都是单根。
\item
  K-代数 \(\mathbf{K}[\mathbf{X}] / (f)\) 是平展的（V，第28页，定义1）。
\end{enumerate}

\(a) \Rightarrow d)\) ：在假设 \(a)\) 下，存在两个多项式 \(g\) 与 \(h\) （在 K{[}X{]} 中），使得 \(fg + f'h = 1\) （IV，第～12页）。设 \(a\) 为 \(f\) 在 \(\Omega\) 中的一个根；我们有

\[
f ^ {\prime} (a) h (a) = f (a) g (a) + f ^ {\prime} (a) h (a) = 1,
\]

由此 \(f'(a) \neq 0\) ；因此 \(a\) 是 \(f\) 在 \(\Omega\) 中的一个单根（IV，第17页，命题7）。 \(d) \Rightarrow c)\) ：若 \(d)\) 成立，则 \(f\) 在 \(\Omega[X]\) 中分解为若干个次数 \(\leqslant 1\) 的不同因子的乘积。

\(c) \Rightarrow b)\) ：在假设 \(c)\) 下，存在 L 中的元素 \(\lambda \neq 0\) 以及 L 的不同元素 \(\alpha_{1}, \ldots, \alpha_{n}\) ，使得 \(f(X) = \lambda(X - \alpha_{1}) \ldots (X - \alpha_{n})\) 。若 \(\deg(f) > 0\) ，我们有（IV，第83页，命题11），

\[
\operatorname{dis} (f) = \lambda^ {2 n - 2} \prod_ {i <   j} (\alpha_ {i} - \alpha_ {j}) ^ {2} \neq 0.
\]

\(b) \Rightarrow a)\) ：设 \(c\) 为首项系数，D 为 \(f\) 的判别式；\(f'\) 与 \(f\) 的结式等于 \(\pm cD\) （IV，第84页，公式(54)），因此非零；因此（IV，第78页，推论2）多项式 \(f\) 与 \(f'\) 在 K{[}X{]} 中互素。

\(a) \Rightarrow e)\) ：设 A 为 K-代数 \(\mathbf{K}[\mathbf{X}] / (f)\) ，并设 \(x\) 为 X 在 A 中的像；由 III，第 573 页，命题 22，A-模 \(\Omega_{\mathrm{K}}(\mathbf{A})\) 由元素 \(dx\) 生成，它们满足唯一的关系 \(f'(x) dx = 0\) 。根据 V，第 33 页，定理 3，K-代数 A 是平展的当且仅当 \(f'(x)\) 是 A 中的可逆元素，而这意味着 \(f\) 与 \(f'\) 在 \(\mathbf{K}[\mathbf{X}]\) 中互素。

定义 2. --- 一个多项式 \(f \in \mathbf{K}[\mathbf{X}]\) 称为可分的，如果它非零且满足命题 3 的等价条件 a)、b)、c)、d) 和 e)。

注记. --- 1) 设 L 是 K 的一个扩张，f 是 K{[}X{]} 中的一个非常数多项式。由命题3的 e) 及 V，第32页，推论2，无论把 f 视为 K{[}X{]} 中的元素还是 L{[}X{]} 中的元素，假设 f 可分都是同一回事。另一方面，完全可能发生这样的情况：f 在 K{[}X{]} 中不可约，但在 L{[}X{]} 中不然。

\begin{enumerate}
\def\labelenumi{\arabic{enumi})}
\setcounter{enumi}{1}
\tightlist
\item
  设 \(f \in \mathbf{K}[\mathbf{X}]\) ；我们知道（IV，第13页，命题13），存在不可约多项式 \(f_1, \ldots, f_m\) （在 \(\mathbf{K}[\mathbf{X}]\) 中），使得 \(f = f_1 \ldots f_m\) 。设 \(\Omega\) 为 \(\mathbf{K}\) 的一个代数闭包；由于不可约多项式 \(g \in \mathbf{K}[\mathbf{X}]\) 是 \(\mathbf{K}\) 上其每个根（在 \(\Omega\) 中）的极小多项式，\(\mathbf{K}[\mathbf{X}]\) 中两个不同的不可约多项式在 \(\Omega\) 中没有公共根。命题3的条件 \(d\) ) 现在表明 \(f\) 是可分的当且仅当多项式 \(f_1, \ldots, f_m\) 是可分的且两两不同。
\end{enumerate}

命题4. --- 设 \(f\) 是 \(\mathbf{K}[\mathbf{X}]\) 中的一个不可约多项式。则以下条件等价：

\begin{enumerate}
\def\labelenumi{\alph{enumi})}
\item
  \(f\) 是可分的。
\item
  存在一个扩张 \(\mathbf{L}\) （\(\mathbf{K}\) 的），\(f\) 在其中有一个单根。
\item
  导数 \(f'\) （\(f\) 的）不为零。
\item
  域 \(\mathbf{K}\) 的特征为 0，或者其特征为 \(p \neq 0\) 且 \(f \notin \mathbf{K}[\mathbf{X}^p]\) 。
\end{enumerate}

我们首先注意到 K{[}X{]} 中的不可约多项式不是常数。显然 a) 蕴含 b)（取 K 的一个代数闭包作为 L）。若 x 是 f 在 K 的扩域 L 中的单根，则有 \(f'(x) \neq 0\) （IV，第 17 页，命题 7），所以 b) 蕴含 c)，而 c) 与 d) 的等价性由 V，第 9 页，推论得出。

最后假设 \(f' \neq 0\) ；设 \(x\) 为 \(f\) 在代数闭扩域 \(\Omega\) （\(K\) 的）中的一个根。由于 \(f\) 是 \(x\) 在 \(K\) 上的极小多项式，且 \(\deg f' < \deg f\) ，我们有 \(f'(x) \neq 0\) ，因此 \(x\) 是 \(f\) 的单根（IV，第 17 页，命题 7）。因此 \(f\) 是可分的，于是我们证明了 \(c)\) 蕴含 \(a)\) 。

推论 1. --- 要使域 K 是完全域，必要且充分的是 K{[}X{]} 的每个不可约多项式都是可分的。

若域 K 的特征为 0，则 K 是完全的，且由上面的 d)，K{[}X{]} 的每个不可约多项式都是可分的。接下来假设 K 的特征为 \(p \neq 0\) 。

首先假设 K 是完全的。我们有 \(K[X^{p}] = K[X]^{p}\) ，因此不存在 K{[}X{]} 的不可约多项式属于 \(K[X^{p}]\) 。根据命题 4，K{[}X{]} 的每个不可约多项式于是都是可分的。

接下来假设 K 是非完全的，因此 \(K \neq K^{p}\) 。设 a 为 K 中不属于 \(K^{p}\) 的一个元素；多项式 \(X^{p} - a\) 在 K{[}X{]} 中不可约（V，第 24 页，引理 1），且它属于 \(K[X^{p}]\) ，因此不是可分的。

推论 2. —— 设 \(f \in \mathbf{K}[\mathbf{X}]\) 是一个非零多项式。要使 \(f\) 可分，必要且充分的是：存在一个扩张 \(L\) （\(K\) 的），它是一个完全域，且 \(f\) 在 \(L[X]\) 中没有重复因子。

设 \(\Omega\) 为 \(\mathbf{K}\) 的一个代数闭包；若 \(f\) 是可分的，则 \(f\) 在 \(\Omega[\mathbf{X}]\) 中没有重复因子（命题3， \(d\) )）。反之，若 \(\mathbf{L}\) 是 \(\mathbf{K}\) 的一个完全扩张，使得 \(f\) 在 \(\mathbf{L}[\mathbf{X}]\) 中没有重复因子，则 \(f\) 在 \(\mathbf{L}[\mathbf{X}]\) 中是可分的（推论1与注记2），因此在 \(\mathbf{K}[\mathbf{X}]\) 中也是可分的（注记1）。

